\begin{chapter}{Conclusão}\label{cap:conclusao}

A avaliação gerou a hipótese, \textit{post hoc} e a validar, de que abaixo do corte do OpenRecLink, região que reúne 61.485 dos 61.696 potenciais pares do extrato, a ordenação por probabilidade ou escore concentra os vínculos e permite reduzir a revisão necessária para uma mesma recuperação, em relação à ordenação pelo escore do próprio OpenRecLink (O1), pela revisão de menos pares na ordem escolhida, não pela política de quatro ações. Na vigilância, O1 requer 501 revisões para recuperar 23 PAR~=~2, contra 125 de O3 e 96 revisões humanas de P, estas acompanhadas de 344 pares avaliados pelo LLM (\path{data/comparador_mesma_carga.csv}). A hipótese seria contrariada se, em dados novos, a ordenação por probabilidade não concentrasse os vínculos no topo da fila ou se a carga necessária para recuperar 23 vínculos na vigilância não fosse menor que a de O1.

No arranjo sequencial P, resultado principal com registro interno, o corte é preservado e a política com recalibração isotônica decide os potenciais pares abaixo dele; O3 é a fila ordenada pela probabilidade isotônica fora da amostra, e O2, a fila ordenada pela nota do CdR (Seções~\ref{sec:e4-arranjo-metodo} e~\ref{sec:e4-exploratorias-metodo}). Com o \texttt{MATCH} do revisor LLM, Claude Haiku 5.5, encaminhado ao humano e revisor perfeito nos dois lados, P usaria 121 revisões para 23 PAR~=~2 na vigilância; O3 recuperaria 22 na mesma carga e igualaria P com 125. Na confirmação, ambos recuperariam nove com 29 revisões. A vantagem é de no máximo um vínculo, ao custo de 344 potenciais pares avaliados pelo LLM na vigilância, e uma diferença de um vínculo em 36 não sustenta superioridade. Sob a mesma hipótese de revisor e na mesma carga, nenhuma domina na vigilância: com 96 revisões humanas, O3 recupera 22 PAR~=~2, P recupera 23 com 344 pares no LLM e 21 FP, e O2 recupera 12. Na confirmação, a ordenação pela nota do CdR domina fracamente P, pois recupera os mesmos nove PAR~=~2 com 11 revisões, sem FP e sem LLM, contra 17 revisões, 41 pares no LLM e 4 FP de P (\path{data/comparador_mesma_carga.csv}).

O LLM não se paga abaixo do corte como componente da política de quatro ações. Na conta exploratória em revisões humanas equivalentes, com revisor perfeito, os custos de cada modo e a parcela dos vínculos perdidos, a perda total na vigilância é de 957,72 para P, 775 para O3 com 125 revisões e 772,72 para P com \texttt{MATCH} ao humano, os três com 23 PAR~=~2 recuperados, e de 1.296 para O2 na carga de P, com 12 recuperados; na confirmação, é de 957,2 para P, 569 para O3, 569,2 para P com \texttt{MATCH} ao humano e 551 para O2, todos com nove (\path{data/custos_equivalentes_exploratorios.csv}). A conta toma como zero o FP humano, embora a política adote proporção de erro humano de 0,05, e não é uma avaliação econômica do filtro exploratório (Seções~\ref{sec:e4-mesma-carga} e~\ref{sec:e4-filtro}) \cite{Jardim2026phdrepo}.

\section{Síntese dos principais achados}\label{sec:concl-sintese}

No primeiro objetivo (Capítulo~\ref{cap:estudo1}), o comparador de registros foi reimplementado em Python, documentado e disponibilizado, e descrito com 29 subescores de similaridade (7 de nome, 7 de nome da mãe, 5 de data, 4 de município e 6 de endereço), dos quais 28, com a nota final, formam as 29 colunas-base que, somadas a 29 atributos derivados, totalizam 58 atributos no vetor de entrada. Pela importância SHAP, os atributos de maior contribuição no XGBoost foram \texttt{NOMEMAE\_qtd\_frag\_comuns} (0,932), \texttt{diff\_ano\_nasc} (0,768) e \texttt{NOME\_qtd\_frag\_muito\_parec} (0,722); 11 das 20 variáveis de maior importância média (55\%) são compartilhadas com Random Forest, com correlação de postos de Spearman de 0,38.

A verificação de equivalência parcial com o protótipo em Java, em lote sintético, mostrou concordância nos subescores de data por igualdade e por um dígito e divergências nos subescores de frequência, no Soundex e nas inversões de data, cujos mecanismos foram identificados por inspeção do código, sem que a comparação eleja uma das implementações como correta, de modo que a equivalência entre elas não foi demonstrada; a nota final coincidiu em 11,2\% dos pares, valor que não estima a concordância nos dados do estudo.

No segundo objetivo (Capítulo~\ref{cap:estudo2}), a ordenação dos classificadores dependeu da métrica: com limiar 0,5, os modelos baseados em árvores tiveram $F_1$ mais elevado que a regressão logística, ajustada com pesos de classe balanceados, mas a área sob a curva de precisão e sensibilidade (AUC-PR) da regressão logística foi de 0,9667, contra 0,9564 do Random Forest (\path{data/nb01_results.json}). Como nenhum experimento variou o grau de desbalanceamento, esses resultados não sustentam superioridade das árvores sob desbalanceamento extremo. No QW1, o XGBoost obteve o maior $F_1$ médio entre os classificadores avaliados, com $0{,}924 \pm 0{,}030$ na validação cruzada $5 \times 5$ realizada na parcela de treinamento. Na análise original, a combinação RF+SMOTE alcançou $0{,}916 \pm 0{,}026$ em cinco partições estratificadas da base completa. São protocolos distintos, que não permitem ordenar essas duas configurações sob as mesmas condições de avaliação. A estratificação por faixas da nota final do CdR mostrou que a maior parte dos erros e ambiguidades se concentra no intervalo intermediário, no qual a combinação de múltiplas evidências (nominais, data de nascimento, endereço e município) é necessária para reduzir a incerteza \cite{Breiman2001rf,Chen2016xgboost,Tyagi2025pprl}.

O terceiro objetivo comparou estratégias operacionais de vigilância e de confirmação quanto aos vínculos recuperados, aos falsos positivos e à carga de revisão na partição de teste (Capítulo~\ref{cap:estudo3}). Na ablação, o $F_1$ foi de 0,931 para Random Forest com SMOTE e de 0,943 para o XGBoost isolado no fluxo Modal, e a melhor combinação desse fluxo, XGBoost com probabilidade $\geq 0{,}7$ ou escore de regras $\geq 8$, alcançou 0,958, com 68 vínculos detectados, nenhum falso positivo e seis perdidos, todos com limiar escolhido no teste (\path{data/ablation_results.csv}). Os indicadores de encerramento por óbito constituem vazamento de alvo nesses experimentos. Por isso, os resultados não fundamentam a adoção dos \textit{pipelines} avaliados nem demonstram ganho epidemiológico dessas estratégias. As 9 estratégias de balanceamento foram avaliadas apenas com Random Forest; seu efeito não se transfere automaticamente ao XGBoost.

Nesse objetivo, as estratégias orientadas à vigilância e à confirmação foram descritas por fronteiras de Pareto com dez pontos não dominados dentro de cada modelo, sete do RF+SMOTE e três do XGBoost (Tabela~\ref{tab:pareto-frontier}). No teste da ablação original, o RF+SMOTE com limiar 0,5 detectou 67 vínculos, contra 43 da decisão hipotética de limiar $\geq 8$ da nota final, que não corresponde ao corte por passo aplicado pelo serviço (\path{data/epidemiological_impact.csv}). Dos 27 vínculos detectados pelo classificador e não por essa decisão, 25 são PAR~=~1, já aceitos pelo corte do OpenRecLink, e dois são PAR~=~2; esses 27 vínculos pertencem a 24 registros de óbito distintos, e os 247 vínculos do extrato, a 204 (\path{data/banca_b_recuperados_por_par.json}). As razões de 70 revisões para 67 vínculos detectados, cerca de 1,0, e de 67 para 43, cerca de 1,6, são o inverso do VPP de cada decisão, não custo por óbito. A decisão pela nota $\geq 8$ teve $F_1$ de 0,610, contra 0,842 do braço somente regras (regras $\geq$ 6) (Tabela~\ref{tab:ablation-best-category}). Sob vazamento de alvo, partição por linha e limiares explorados no teste, esses valores descrevem a comparação interna, não a redução da subnotificação.

O quarto objetivo avaliou o GZ-CMD abaixo do corte, preservando os 211 vínculos definidos por ele, frente à referência operacional produzida pelo corte e pela revisão do serviço e, quanto à carga, frente a ordenações de mesma carga. Na vigilância, P aceita automaticamente 215 dos 247 vínculos e encaminha 96 potenciais pares ao revisor humano (\path{data/arranjo_corte_metricas.csv}). P não reproduz 13 e 27 dos 36 vínculos recuperados pelo serviço, perdidos por decisão de não vínculo sem chegar ao revisor; aceita 25 e 12 potenciais pares abaixo do corte, dos quais 21 e 4 foram rejeitados ou não examinados pelo serviço e 4 e 8 são PAR~=~2. Na vigilância, 17 dos 21 FP situam-se acima do limiar global hipotético e 11 acima do limiar por passo; na confirmação, os 4 estão acima de ambos (Seção~\ref{sec:e4-mesma-carga}). A escolha depende do orçamento: na confirmação, O2 recupera 9 vínculos com 11 revisões, enquanto O3 requer 29; para recuperar os 36, O1 requer 2.521, posição do último vínculo nessa fila, e O3, 23.443. A 20 minutos por par, as 96 revisões de P ou de O3 na vigilância corresponderiam a 32 horas (\path{data/esforco_arranjo.csv}); a redução pertinente é frente a O1 na mesma recuperação, que requer 501 revisões para 23 vínculos, e não frente ao cenário condicional de 2.156 revisões, de meta diferente e que não mede carga histórica (Seção~\ref{sec:padrao-ouro}). Os cenários mistos humano/LLM da Tabela~\ref{tab:cap8-esforco} não são equivalentes a horas humanas observadas. O corte único da nota em 9 recuperaria 37,7\%; a revisão integral dos 21.620 pares da área cinza da nota projetaria 95,7\%, sob sensibilidade humana hipotética de 0,95 (Tabela~\ref{tab:cap8-referencia}).

As análises exploratórias \textit{post hoc} do quarto objetivo acrescentam dois resultados. Como filtro de \texttt{NONMATCH} antes do revisor humano, com 96 pares enviados ao LLM na vigilância, restaram 22 revisões humanas, com 13 PAR~=~2 recuperados, contra seis de O3 na mesma carga humana; o filtro descartou, porém, 9 dos 22 PAR~=~2 do topo, com sensibilidade de 59,1\% (13/22) e VPN de 87,8\% (IC 95\%: 78,5 a 93,5), e esses vínculos, pela Proposição~\ref{prop:duas-acoes-orcamento}, deveriam ser revisados (\path{data/filtro_llm.csv}). A vantagem existe somente com orçamento humano fixo em 22 revisões: na unidade de custo da tese, a perda do filtro, de 1.172,5, supera a de O3 com 96 revisões, de 796 (Seção~\ref{sec:e4-filtro}; \path{data/custos_equivalentes_exploratorios.csv}). No treinamento, a subamostragem dos negativos nas razões 1:1 e 1:5 degradou o topo da fila, com mediana de PAR~=~2 entre os 96 primeiros de 13 e 18,5, contra 22 da referência; as razões 1:20 e 1:50 foram indistinguíveis da referência, comparação limitada porque a referência teve uma execução e cada razão, dez sementes. A mediana do escore de Brier e do ECE da isotônica foi pior em todas as razões, por margem mínima em 1:50 (\path{data/subamostragem_treino_resumo.csv}).

Os valores de destaque do exame de qualificação, $F_1$ de 0,954 na vigilância e precisão de 0,957 na confirmação, provêm da ablação do GZ-CMD com a política v3, cinco rodadas de partição por \textit{comprec}, \textit{Platt scaling} da nota ajustado na partição de treinamento e revisão assistida pelo Kimi K2.5 (Tabelas~\ref{tab:ablacao-comprec-vigilancia-macd-on} e~\ref{tab:ablacao-comprec-confirmacao-macd-on}); essa ablação não foi reavaliada com a política v3.1 nem com o dossiê sem as saídas do motor. Na mesma política histórica, a cascata v3 sobre a base inteira perdeu 91 vínculos por \texttt{NONMATCH} na vigilância, com sensibilidade automática de 58,3\% (Tabela~\ref{tab:cap8-referencia}). O arranjo P, com a política v3.1, a calibração isotônica fora da amostra agrupada por REFREC e Claude Haiku 5.5 com dossiê sem as saídas do motor, teve sensibilidade automática de 87,0\% e 88,7\% e VPP de 91,1\% e 98,2\% (\path{data/arranjo_corte_metricas.csv}), valores que não são comparáveis aos da ablação. A redução da revisão clerical de 21.620 para 41 pares, de 99,8\%, contrapunha a área cinza da nota aos 41 inconclusivos encaminhados ao humano pela cascata com Kimi K2.5, no lote histórico de 1.410 pares selecionado pela política v3, com dossiê que incluía as saídas do motor e probabilidade de \textit{Platt scaling} ajustada nos próprios \textit{labels} (\path{data/exp2_reruns_resumo.csv}). Com a política v3.1, a carga correspondente é de 96 e 17 encaminhamentos humanos no arranjo e de 111 e 127 na cascata v3.1 sobre a base inteira, isto é, aplicada aos 61.696 potenciais pares sem preservar previamente as aceitações do corte e mantida como comparação, trabalho pendente e não acerto observado. Os 18 minutos correspondem à duração paralela desse lote, que dominava o processamento computacional; não medem o processo de \textit{linkage} nem o esforço humano (1.066,8\,s com 15 trabalhadores e mediana de 9,7\,s por par; \path{data/exp2_latencias_revisores.csv}).

No quinto objetivo (Capítulo~\ref{cap:estudo5}), a disparidade por passo do arranjo é herdada do corte e do \textit{label}: o corte aceita 15 de 40 vínculos no passo~11 (37,5\%), contra 95,0\% e 94,5\% nos passos~6 e~9. Dos 25 PAR~=~2 do passo~11, o arranjo aceita dois na vigilância e quatro na confirmação \footnote{Fonte: \path{data/par_por_passo.csv}.}. Sua sensibilidade é de 42,5\% e 47,5\%, contra pelo menos 95,0\% nos passos~6 e~9; nessa análise, o falso negativo inclui todo vínculo não aceito automaticamente, inclusive os encaminhados ao humano, sentido mais amplo que o dos vínculos perdidos por não vínculo no quarto objetivo. Na vigilância, a sensibilidade no passo~11 fica abaixo dos 45,0\% da cascata sobre a base inteira. Houve três rejeições em 26 comparações na vigilância e duas em 26 na confirmação após Benjamini-Hochberg, contra seis em 28 e uma em 21 nessa cascata, todas operacionais. Nos testes pareados, consideraram-se somente as comparações em que os dois grupos têm ao menos 30 observações no denominador da métrica. Todas as comparações entre quartis do nome foram excluídas pelo filtro, assim como as de município e de completude, de modo que, nessas dimensões, zero rejeições significa ausência de comparação elegível, e não paridade. Ausência de rejeição demográfica não demonstra equidade. O VPP do arranjo na vigilância foi de 86,3\% nas mulheres e 93,3\% nos homens, com proporção de falsos positivos entre os não vínculos (FPR, \textit{false positive rate}) de 0,078\% e 0,023\%. Os VPP de 93,1\% e 70,7\% e as FPR de 0,04\% e 0,14\% pertencem ao XGBoost; os 7.452 positivos, ao corte 6,0 da nota. A incompatibilidade entre paridade de sensibilidade e VPP sob prevalências desiguais \cite{Kleinberg2016inherent,Chouldechova2017fair} não dispensa analisar as proporções de erro; a ausência de raça/cor impede avaliar esse eixo.

\section{Contribuições}\label{sec:concl-contribuicoes}

As contribuições compreendem a decomposição por PAR, que distingue a preservação dos vínculos definidos pelo corte do ganho sobre os recuperados pela revisão; a curva de esforço aplicada a este fluxo como instrumento de escolha da ordenação e do ponto de parada, sem reivindicar originalidade da ordenação por pesos ou da revisão clerical descritas por \citeonline{Sariyar2010recordlinkage}; e a investigação do LLM como aceitador e como filtro de descarte antes do revisor humano, associada à investigação de sua ancoragem. Nesta última, o VPP de 16,0\% e os 9,06 FP previstos contra 21 observados são resultados \textit{post hoc}; como filtro de \texttt{NONMATCH}, o mesmo revisor apresentou VPN de 87,8\% e sensibilidade de 59,1\% na vigilância, em $k=96$; em outro orçamento, $k=125$, descartou 10 dos 23 vínculos. A utilidade do filtro é exploratória, não validação de sua adoção. Completam as contribuições a política auditável, o CdR em Python, o pacote de código e os resultados agregados públicos, o \textit{benchmark} de acesso restrito e o arranjo P com registro interno no repositório antes da avaliação com os \textit{labels} (registro \texttt{8c16a41}) \cite{Jardim2026phdrepo}. Ler esse arranjo como resultado negativo é interpretação posterior: o registro listou desfechos sem hipótese direcional nem critério de sucesso, a referência recupera os 36 PAR~=~2 por definição do \textit{label}, e o código inicialmente registrado não está no histórico do repositório (Seções~\ref{sec:e4-arranjo-metodo}, \ref{sec:e4-mesma-carga} e~\ref{sec:e4-filtro}). A referência interna de P não lhe confere força confirmatória. O \textit{benchmark} e o pacote apoiam reprodução, não validação externa \cite{Jardim2026gzcmd,Jardim2026zenodo}.

Os artefatos do GZ-CMD v3.1, compreendendo funis por potencial par, probabilidades OOF agrupadas por REFREC, respostas dos revisores identificadas por \texttt{pair\_id}, proporções de erro por banda e instruções de revisão (\textit{prompts}) v3 e v4, integram a versão 3.0.0 do depósito Zenodo, utilizada nos experimentos (DOI 10.5281/zenodo.23928226) \cite{Jardim2026zenodoGzcmd}, com as classificações da referência operacional por potencial par, sem identificadores diretos. A versão corrente, 3.0.1, corrige documentação e metadados, mantendo os resultados experimentais \cite{Jardim2026zenodoGzcmd301}; a licença MIT se aplica somente ao código. O depósito é público com arquivos restritos, mediante solicitação e acordo de uso, sob os pareceres do CEP descritos na Seção~\ref{sec:dados-acesso-etica}, e inclui, além dos escores, a base de análise de equidade por potencial par, com faixa etária, sexo, unidade da federação e indicador de encerramento por óbito; os \textit{labels} e essa base por potencial par constam somente do depósito restrito. Os dados permanecem pseudonimizados para o laboratório, e o \texttt{pair\_id}, calculado sem sal, não constitui barreira criptográfica suficiente contra a enumeração de identificadores internos. O conteúdo depositado e os limites das medidas de k-anonimidade constam da Seção~\ref{sec:ambiente-computacional}.

A exportação restrita com pseudônimo por HMAC-SHA256, apresentada na Seção~\ref{sec:export-restrita}, foi implementada para novas exportações, sem substituir os identificadores dos artefatos depositados; a nova matriz não foi publicada.

O depósito Zenodo contém a política v3.1 sobre a base inteira, não o arranjo sequencial. Os agregados do arranjo estão no repositório da tese; as saídas por potencial par permanecem locais \cite{Jardim2026phdrepo}.

O GZ-CMD documenta o que deixa de recuperar. O contraste por PAR com o comparador sem recalibração está na Tabela~\ref{tab:cap8-arranjo-corte-par2}. A distinção entre aceitação automática, encaminhamento pendente e cenário com revisão humana impede interpretar redução de carga como manutenção dos vínculos do serviço. P é o braço registrado, e a recalibração isotônica, cujo conjunto de calibração na primeira partição externa tinha 32 PAR~=~2, é decisão de registro, não superioridade demonstrada. P-Platt constava como braço no registro e passou a sensibilidade na terceira emenda, e o critério de estabilidade de 18/20 foi fixado depois da avaliação original (Seção~\ref{sec:e4-arranjo-metodo}; \path{gzcmd/M_arranjo_corte_desvios.md}). Pela função de perda declarada, aplicada às medianas das réplicas, P-Platt teve perda menor que P nos quatro cenários, com diferenças de 144,1 e 94,4 na vigilância e de 101,6 e 143,0 na confirmação, nos blocos numérico e de partição (\path{data/robustez_perda_ponderada.csv}). A dominância pontual de P-Platt sobre P não se sustentou entre réplicas: ocorreu em 7/20 e 4/20 réplicas dos blocos numérico e de partição na vigilância e em 7/20 e 5/20 na confirmação, com dominância inversa em 0/20. P-Platt teve menos FP que P em 20/20 réplicas na vigilância, e a probabilidade isotônica oscilou mais, com variação máxima mediana de 0,46 contra 0,17 no bloco numérico (Seção~\ref{sec:e4-robustez}; \path{data/robustez_arranjo_resumo.csv} e \path{data/robustez_perda_ponderada.csv}). A estabilidade não sustenta, portanto, superioridade de P, que permanece resultado principal por ter sido o braço registrado. A análise de estabilidade descreve a estabilidade interna entre réplicas, não validação em novos dados \cite{Jardim2026phdrepo}.

\section{Limitações e generalização}\label{sec:concl-limitacoes}

A principal limitação reside na referência operacional e na seleção dos casos verificados. A epidemiologista que executou o \textit{linkage} de origem e o CdR foi a revisora única, sem classificação independente. Para os passos 6, 9 e 11, a ordem entre o CdR e a revisão e o que a revisora consultou não estão documentados; adota-se como hipótese de trabalho que o CdR precedeu a revisão, com nota e subescores disponíveis à revisora, o que configuraria viés de verificação em direção a O2 a O4, de magnitude não estimada (Apêndice~\ref{ap:linkage-origem}). O pré-processamento atribuído a outro profissional não constitui segunda avaliação. Se a revisão tivesse seguido o escore do OpenRecLink, os PAR~=~2 recuperados tenderiam a se concentrar no topo dessa ordem, o que favoreceria O1, mas a ordem não tem fonte. Os PAR~=~1 ficam fora da comparação abaixo do corte, e O3 e O4 também foram treinadas sobre os mesmos \textit{labels} verificados. A possível incorporação das evidências do CdR à referência, com potencial viés de incorporação, não foi quantificada (Seção~\ref{sec:padrao-ouro}). Revisão cega e independente em dados novos permitiria investigar esses vieses \cite{Harron2017linkagequality,Oliveira2026obitoTBpulmonar}.

Uma segunda limitação refere-se ao risco de falsos negativos na referência operacional decorrentes de falhas na etapa de blocagem (\textit{blocking}). Se um vínculo não foi gerado como candidato pelo OpenRecLink em nenhum dos passos de blocagem, ele estará ausente do universo avaliado, e tanto o classificador quanto a revisão manual serão incapazes de recuperá-lo. Essa limitação é inerente a qualquer estudo de \textit{linkage} probabilístico baseado em blocagem \cite{Doidge2019linkageerror} e implica que as métricas de sensibilidade reportadas nesta tese representam estimativas condicionais ao conjunto de candidatos gerados, não estimativas absolutas da capacidade de detecção. Abordagens alternativas, como a geração de identidades sintéticas \cite{Lam2024synthetic} ou a validação cruzada com múltiplas estratégias de blocagem, poderiam mitigar parcialmente essa restrição.

A sensibilidade da blocagem, no que se refere aos vínculos verdadeiros nunca submetidos à comparação, não é estimável com o extrato atual. A auditoria M8 registrou zero recapturas entre os 247 vínculos distintos dos três passos exportados (Seção~\ref{sec:blocagem-rendimento}), resultado compatível com a exclusão progressiva dos registros já pareados pelo OpenRecLink, que impede que os pares formados nos passos anteriores sejam novamente observados nos seguintes. Essa ausência de sobreposição entre vínculos decorre do desenho sequencial e inviabiliza a estimação por captura-recaptura entre passos, não constituindo evidência de completude da blocagem. Sua investigação exigiria nova execução e exportação sem exclusão entre passos, sobre o mesmo universo de origem e com preservação do histórico de cada par, além da avaliação das hipóteses necessárias à estimação.

O desbalanceamento extremo (1:249) entre vínculos e não pares constitui desafio estatístico relevante, pois pequenas variações na proporção de falsos positivos podem gerar grandes volumes de revisão. A análise de sensibilidade com nove estratégias de reamostragem registrou F$_1$-Score de 0,880 a 0,918, descrevendo variação interna sem demonstrar benefício uniforme dos tratamentos. A transposição direta dos limiares selecionados para bases com proporções de desbalanceamento distintas requer recalibração, uma vez que os valores preditivos positivo e negativo dependem da prevalência \cite{Shaw2022biases}.

A disponibilidade restrita de variáveis clínicas no conjunto exportado pelo comparador de registros limita análises epidemiológicas mais detalhadas. Raça/cor e comorbidades, embora presentes nas bases originais do SIM e do Sinan, não constam do extrato; o bairro de residência consta dele, mas não integrou o vetor de características, por razões de escopo e privacidade, e foi usado somente na completude da análise de equidade. A análise de equidade apresentada na Seção~\ref{sec:e5-resultados-equidade} contornou parcialmente essa restrição ao estratificar as decisões em subgrupos demográficos (sexo e faixa etária) e operacionais (quartil de concordância do nome, situação de encerramento da notificação, município, passo de blocagem, ano do óbito e completude). Essa análise mostrou disparidade dominante no passo de blocagem. A direção da diferença de VPP por sexo no XGBoost se inverte no arranjo; não se deve transferir a interpretação de uma decisão para a outra. A equidade do processo completo depende ainda da revisão humana dos encaminhamentos, cuja qualidade não foi observada. Eixos de estratificação reconhecidamente associados a disparidades em saúde, como raça/cor, ausente do extrato, e geolocalização mais fina que o bairro, permanecem inacessíveis nesse conjunto de dados, o que configura lacuna a ser endereçada em investigações futuras \cite{Mehrabi2021fairness}. A literatura documenta que a qualidade do preenchimento varia por região e por sistema de informação \cite{Lima2020tbquality, Bartholomay2020drtb}; portanto, a generalização para outras localidades e pares de bases (por exemplo, SIM$\times$SIH ou Sinan$\times$GAL) depende de revalidação externa com referências locais \cite{Coeli2021suboptimal}.

A avaliação de arquiteturas de aprendizado profundo constitui questão recorrente na literatura recente de \textit{linkage}. Abordagens fim-a-fim como DeepMatcher \cite{Mudgal2018deepmatcher} e DITTO \cite{Li2020ditto} recebem campos textuais brutos e aprendem representações via modelos de linguagem pré-treinados, reportando ganhos de 1 a 5 pontos percentuais de $F_1$ sobre métodos tradicionais em \textit{benchmarks} com milhares de exemplos rotulados. Neste estudo, o vetor de entrada é tabular, com 29 colunas-base do comparador, incluindo a nota final, e 29 atributos derivados. A literatura oferece antecedentes para a comparação entre classificadores arbóreos e neurais nesse tipo de representação \cite{Grinsztajn2022tabular,ShwartzZiv2022tabular}. A indisponibilidade de campos textuais brutos restringe a avaliação de arquiteturas fim-a-fim que operam diretamente sobre nomes e endereços. Os 247 vínculos delimitam a amostra positiva disponível, mas não demonstram, isoladamente, insuficiência para ajuste fino nem explicam causalmente o desempenho neural. O QW1 incluiu duas MLP e o TabNet \cite{Arik2021tabnet}; na validação cruzada $5 \times 5$, o TabNet apresentou desempenho inferior aos classificadores arbóreos avaliados (Seção~\ref{sec:qw1-deep-learning}), sem que esse resultado se estenda a outras arquiteturas ou configurações.

No GZ-CMD, a seleção para revisão assistida decorre da comparação de \textit{expected loss}, não de demonstração de que o classificador tabular esgotou seu poder discriminativo. O viés do revisor LLM descrito no Estudo 4 limita a interpretação dos resultados históricos, e o desempenho sem as saídas do motor precisa ser avaliado com instruções compatíveis com o dossiê. A avaliação de arquiteturas fim-a-fim com acesso controlado a campos textuais requer estudo próprio, sem pressupor ganho sobre a representação numérica atual \cite{Jardim2026phdrepo}.

A calibração absoluta, a ancoragem do revisor e a ausência de validação externa delimitam a generalização do GZ-CMD (Seção~\ref{sec:e4-limitacoes}). PAR~=~1 é definido pelo corte, com zero FP acima dele por construção. A reutilização das coortes e as chamadas novas estão descritas na Seção~\ref{sec:e4-arranjo-metodo}; a avaliação não é prospectiva. O enquadramento sequencial foi escolhido \textit{a posteriori}. O registro interno antecedeu a S3, mas não essa escolha, e o código inicialmente registrado não está disponível no histórico; os desvios estão documentados no Estudo~4. As proporções do revisor, estimadas com denominadores pequenos e pares mais fáceis, não se transferiram como previsto à região abaixo do corte. Nem P-Platt nem o comparador sem recalibração podem tornar-se resultado principal sem novo registro. Os indicadores de encerramento por óbito integram QW1, ablação e comparação com a nota $\geq 8$, constituindo vazamento de alvo ausente do GZ-CMD (Seções~\ref{sec:e2-limitacoes} e~\ref{sec:e3-limitacoes}). A sensibilidade à retirada desses indicadores e à partição por REFREC consta do Apêndice~\ref{ap:sensibilidade}; a ablação das medidas contínuas permanece sem essa reavaliação.

A revisão assistida acrescenta limitações próprias. A geração do Claude Haiku 5.5 é estocástica: na ablação do dossiê, controle e execução original concordaram em 13 de 18 pares; na análise de estabilidade, as respostas do LLM ficaram fixas entre réplicas, de modo que essa variação não entra na estabilidade medida. Dez modelos foram avaliados em lotes das mesmas coortes, e a escolha do revisor de referência e das instruções v4 é seleção \textit{post hoc}, sujeita à maldição do vencedor (Seção~\ref{sec:e4-limitacoes}). Os dossiês pseudonimizados foram enviados a provedores sediados no exterior sem autorização específica dos pareceres para essa transferência internacional, limitação ética descrita na Seção~\ref{sec:dados-acesso-etica}. Toda a avaliação ocorreu nos mesmos dados, com predições fora da amostra apenas internas; nenhuma validação externa foi realizada.

Nos cenários de sensibilidade humana de 0,95, 0,99 e 1,00, a amplitude máxima da sensibilidade total foi de 0,73 ponto percentual na revisão integral e de 0,53 nas estratégias de carga reduzida, sem inversão de ordem entre estratégias com a mesma carga humana (Seção~\ref{sec:e4-sensibilidade-humana}). Esses limites decorrem da construção dos cenários, pois 5\% de erro sobre os 36 PAR~=~2 abaixo do corte, entre 247 vínculos, não ultrapassa 0,73 ponto percentual, e não constituem achado empírico. As perdas são de cobertura, não de erro humano: os vínculos não recuperados não chegam ao revisor. A sensibilidade humana real não foi medida; os 5\% de erro vêm da configuração, e a especificidade humana é tratada como perfeita, sem modelar FP humanos, inclusive na conta de custos, embora a política adote proporção de erro humano de 0,05. Essas hipóteses não eliminam os FP automáticos observados de P.

\section{Recomendações operacionais}\label{sec:concl-recomendacoes}

Recomenda-se, como hipótese condicionada à validação, preservar o corte e, abaixo dele, revisar em ordem decrescente de uma probabilidade ou escore, escolhendo a ordenação e o ponto de parada pela curva de esforço conforme o orçamento humano. Nenhuma ordenação domina em todas as cargas, e a recomendação não fixa O3; na confirmação, a ordenação pela nota do CdR recuperou os mesmos vínculos de P com menos revisões. Se a ordenação usar probabilidade aprendida, a subamostragem dos negativos deve evitar as razões 1:1 e 1:5, que degradaram o topo da fila; as razões 1:20 e 1:50 foram indistinguíveis da referência, em comparação de uma execução contra dez sementes por razão, e a subamostragem não melhorou conjuntamente a região operacional e a calibração de forma estável (Seção~\ref{sec:e4-subamostragem}). O LLM pode atuar somente como filtro de \texttt{NONMATCH} antes do revisor humano, hipótese exploratória a validar em dados novos e restrita a orçamento humano fixo, já que, entre os 96 primeiros pares da vigilância, o filtro descartou 9 dos 22 vínculos, que deveriam ser revisados, e, na unidade de custo da tese, perde para O3 com 96 revisões (Seção~\ref{sec:e4-filtro}); nunca deve aceitar vínculos sozinho, pois, como aceitador abaixo do corte, teve VPP \textit{post hoc} de 16,0\%, com quatro vínculos em 25 aceitações na vigilância (\path{data/arranjo_corte_vpp_llm.csv}). Se for utilizado, seu \texttt{MATCH} vai ao revisor humano. Os potenciais pares fora do orçamento recebem não vínculo, decisão do GZ-CMD restrito às ações revisor humano e não vínculo. A adoção depende de validação registrada em dados novos, não de escolha retrospectiva de uma fila favorável \cite{Jardim2026phdrepo}.

As estratégias do Estudo~3 requerem nova avaliação sem os indicadores de encerramento por óbito, com limiar escolhido fora do teste e partições agrupadas por registro. Não se recomenda sua adoção com base nos resultados afetados por vazamento. A importância do nome da mãe na análise SHAP sustenta examinar a qualidade desse campo, mas não dispensa essa reavaliação.

\section{Trabalhos futuros}\label{sec:concl-trabalhos-futuros}

A análise do Apêndice~\ref{ap:sensibilidade} delimita a reexecução já realizada; permanecem como continuidade a validação cruzada agrupada dos classificadores do QW1 e a extensão à ablação de medidas contínuas, com seleção de limiares em validação separada e avaliação final agrupada por REFREC.

Como continuidade, destacam-se: (i) validação externa em outros pares de bases e contextos epidemiológicos; (ii) incorporação de variáveis clínicas e temporais adicionais para ampliar análises de impacto; (iii) estratégias de amostragem adaptativa e de aprendizado ativo (\textit{active learning}, em que o modelo indica quais pares devem ser rotulados) para reduzir a dependência de rotulagem manual; (iv) integração do pós-processamento em \textit{pipelines} operacionais com monitoramento contínuo de desempenho e auditoria; (v) avaliação do efeito de diferentes esquemas de blocagem e comparadores sobre a distribuição da área cinza e a estabilidade dos modelos; e (vi) investigação de arquiteturas de aprendizado profundo fim-a-fim \cite{Mudgal2018deepmatcher,Li2020ditto} em bases administrativas de maior volume, com acesso a campos textuais brutos, cenário no qual a disponibilidade de conjuntos positivos da ordem de milhares pode viabilizar o ajuste fino de modelos pré-treinados.
Quanto ao GZ-CMD, a replicação em cenários com múltiplos revisores, em outros municípios e períodos e em pares de bases com prevalências e completudes distintas, bem como a avaliação da sensibilidade a variações na qualidade de preenchimento, constituem passos necessários antes da adoção em escala.

As análises do quarto objetivo abaixo do corte permanecem como validação pendente, não como resultados confirmados: o arranjo sequencial e a ordenação abaixo do corte requerem nova avaliação registrada em outra coorte; a instabilidade da seleção entre réplicas requer exame em dados novos; a escolha da ordenação e do ponto de parada pelo comparador de mesma carga e pela curva de esforço depende de revisão humana observada; o filtro de \texttt{NONMATCH} requer respostas de uma única execução, com protocolo registrado, seguidas de revisão humana efetiva; a escolha entre o treinamento sem subamostragem e as razões 1:20 e 1:50, indistinguíveis nesta coorte, precisa ser examinada fora dela, com repetições também da referência; e os cenários de sensibilidade humana devem dar lugar a dupla revisão com concordância medida e registro dos FP humanos.

A continuidade requer âncoras internas verificadas em amostra independente, associadas a curva monotônica regularizada, em vez de presumir que momentos extremos determinem a calibração intermediária. Além dessas validações, avaliar a transferência da política requer definir orçamento humano compatível com a rotina do serviço. O pacote Python e a aplicação GZ-CMD Sandbox oferecem implementações públicas para reprodução e experimentação, sem substituir essa validação empírica \cite{Jardim2026gzcmd,Jardim2026gzcmdweb,Jardim2026phdrepo}.

Os exercícios M7 e M8, documentados em \path{gzcmd/M7_weight_space.md} e \path{gzcmd/M8_blocagem.md} e não no corpo da tese, exploraram a auditoria de sistemas de \textit{linkage} por modelos de linguagem a partir de parâmetros e produtos agregados, sem dados individuais, por analogia metodológica com a leitura de pesos de redes neurais discutida por \citeonline{Takehi2026weightspace}, em \textit{post} de blog não revisado por pares, sem replicar ou validar esse trabalho. No M8, retomado na Seção~\ref{sec:blocagem-rendimento}, as chaves de blocagem inferidas dos agregados não recuperaram os componentes fonéticos e o ano de nascimento da documentação e permanecem hipóteses a validar contra a configuração efetiva do OpenRecLink; as chaves alternativas K1 (data de nascimento exata, município e sexo) e K4 (unidade federativa no lugar do município) tiveram sensibilidades descritivas de 42,1\% (104/247) e 46,2\% (114/247) dos vínculos conhecidos, com 0,19\% (116/61.696) e 0,21\% (128/61.696) das linhas. Obtidas entre pares já gerados, essas proporções não medem vínculos novos, sensibilidade absoluta ou custo nas bases completas; sua validação exigiria novo \textit{linkage} de origem, executado pelo laboratório nos dados identificados, em ambiente de acesso controlado. No M7, na área cinza da nota, um escore programado obteve área sob a curva ROC de 0,975, frente a 0,873 da nota final e 0,995 do XGBoost com probabilidades OOF agrupadas por REFREC, mantido fixo na reavaliação com treino e teste agrupados por REFREC nos níveis externo e interno, calibração por âncoras restrita aos grupos de treinamento de cada \textit{fold} externo e sem novas chamadas LLM. Nenhum dos exercícios alterou o motor de decisão; ambos geram hipóteses para validação futura.

A estimação do custo implícito do vínculo, da revisão e das perdas permanece como trabalho futuro.

\end{chapter}
